\chapter{Detailed Implementation Guidance}
\label{app:impl}

A useful estimator must connect a reported quantity to the observations that constrain it.
This appendix develops a candidate measurement chain, with explicit assumptions and checks; it does not report a validated prototype.
Coordinate conventions and reference points follow \cref{ch:params}, and hardware limits follow \cref{app:hardware}.
Measured sensor channels, derived physical quantities and estimates dominated by priors must remain distinguishable in the output.
A plausible number, a small filter covariance or agreement with another launch monitor is insufficient evidence on its own.

\section{Radar Signal Chain (OPS243-A)}
\label{app:radardsp}

\subsection{Doppler Scaling and Resolution Budget}

For the illustrative carrier $f_c=\SI{24.125}{\giga\hertz}$, \cref{eq:doppler} gives approximately \SI{71.95}{\hertz} per mph of radial speed.
With uniformly spaced complex I/Q samples at $f_s=\SI{30}{\kilo\hertz}$, the ideal principal frequency interval is $[-f_s/2,f_s/2)$, corresponding to a radial-speed half-width of about 208.48\,mph.
This sampling calculation does not certify the analog passband, firmware, detection sensitivity or total-speed envelope of an OPS243-A installation~\cite{ops243brief,an029}.
At 200\,mph radial speed, the carrier Doppler is already about \SI{14.39}{\kilo\hertz}; rotational sidebands and acceleration can consume the remaining frequency margin.
Confirm sample cadence, I/Q ordering, mixer sign and missing-frame behavior with the actual acquisition mode.

A 128-sample record lasts \SI{4.267}{\milli\second} and has raw Fourier spacing $f_s/128=\SI{234.375}{\hertz}$, equivalent to approximately 3.26\,mph.
Padding it to 4096 points samples the same finite-record spectrum more densely; it does not create independent observations.
Distinguish bin spacing, separation of neighboring components, single-tone estimation precision and absolute speed accuracy.
All depend differently on window shape, signal-to-noise ratio, calibration and model mismatch.

At 3,000\,rpm, mechanical rotation frequency is \SI{50}{\hertz}.
An observation of \SI{20}{\milli\second} gives a \SI{50}{\hertz} unpadded grid; \SI{40}{\milli\second} gives \SI{25}{\hertz}.
Neither duration guarantees that weak spin sidebands can be separated from a strong drifting carrier: sidelobes, missing harmonic orders and frequency variation also matter.
Choose a window and acceptance criterion using injected signals and recorded shots over the intended operating envelope, rather than treating two rotation periods as a universal requirement.

\subsection{Candidate Spin-Rate Estimation}

US8845442 describes candidate spectral traces, harmonic relationships and order assignment~\cite{us8845442}.
The following is an engineering proposal informed by that disclosure, not its exact algorithm or a claim about current commercial firmware:
\begin{enumerate}\itemsep2pt
\item Preserve coherent raw I/Q with acquisition timestamps and configuration metadata.
Associate a post-impact ridge with the ball using temporal continuity and independently justified gates; the strongest peak need not be the ball.
Exclude saturation and unresolved aliasing before interpreting peak locations.
\item Estimate carrier frequency as a continuous function of time.
For a convention with carrier phase $+2\pi\int f_d\,dt$, coherent demodulation uses
\begin{equation}
\label{eq:impldechirp}
y(t)=z(t)\exp\!\left[-j2\pi\int_{t_*}^{t}\widehat f_d(\tau)\,d\tau\right].
\end{equation}
Here $z$ is the recorded complex signal and $t_*$ is a fixed phase origin.
An independent frequency shift per short-time Fourier frame can recenter its peak but does not generally remove intraframe chirp; discontinuous phase corrections also introduce artifacts.
Carrier-estimation error belongs in the residual uncertainty.
\item Search for harmonic families using a declared spectral or time-domain model.
Cepstral or harmonic-product scores can propose candidates, but missing orders, clutter and approximately symmetric patterns can cause octave errors or competing fundamentals.
Test the candidate order assignments and residuals, rather than declaring either method intrinsically robust to missing harmonics.
\item Require prespecified evidence across usable observations and check against an independent spin reference.
Overlapping Fourier windows are correlated, so three consecutive frames are not three independent confirmations.
Report ambiguity, detection failures and the fraction of attempted shots yielding an accepted measurement.
If only a club-conditioned prior is available, label the result estimated and retain the missing-measurement flag.
\end{enumerate}

Scattering contrast, viewing direction and ball pattern influence which modulation components are observable (\cref{sec:radarspin}).
Log ball identity and markings rather than assigning a universal ranking to clean urethane, range or marked balls.
A candidate rotation frequency does not by itself identify the full spin vector.
US8845442 and US9868044 offer distinct technical approaches to study; the latter's dielectric-lens explanation and claimed processing requirements must remain separate~\cite{us8845442,us9868044}.
These disclosures establish neither implementation clearance nor technique-wide expiry dates (\cref{sec:fto}).

\subsection{Club-Speed Extraction}

The velocity of a rigid-head point $P$ is $\vect v_P=\vect v_O+\vect\omega\times\vect r_{OP}$; radar observes its line-of-sight projection, weighted by the actual scattering response (\cref{sec:radarclub,app:screw}).
The heel need not always be the slower return, and shaft returns need not occupy a separable bottom decile.
Trimming the highest and lowest spectral deciles and taking a median defines a reproducible spectral statistic, not a geometric-center speed.
It can become a calibrated proxy only after reference-point, timing and pose dependence have been characterized against an independent measurement.

TrackMan's documented geometric-center speed definition specifies the intended reported quantity, not a reconstruction algorithm that can be assumed for a different sensor~\cite{trackmanclubspeed}.
For a candidate club track, retain the pre-impact observation interval, target-association evidence and uncertainty in projection and reference-point conversion.
A smash-factor discrepancy is a diagnostic lead, not proof of toe-lock or grounds for a universal rejection threshold (\cref{sec:smash}).
Validate the speed estimate before using it to validate an impact model; otherwise the same assumption can appear to confirm itself.

\Needspace{8\baselineskip}
\section{Angle Measurement (K-LD7) and the EKF}
\label{app:ekf}

\subsection{Interferometric Angles and Hardware Limits}

For the declared ideal two-receiver geometry, phase difference obeys \cref{eq:interferometry}.
Phase calibration, wrapping, multiple targets and the antenna response affect the inverse angle estimate.
The K-LD7 revision-B datasheet specifies a 6.223\,mm receiver separation, FSK ranging, processed detections and a raw-data request interface; it does not justify treating the module as an arbitrary FMCW range--Doppler instrument~\cite{kld7}.
Its maximum speed setting is 100\,km/h, with a typical frame duration of \SI{29}{\milli\second} at that setting.
A ball whose radial component exceeds that limit is outside the stated unambiguous speed range; combining sensors in a filter does not make an invalid channel valid.
Thus a \SI{50}{\milli\second} interval cannot be assumed to contain ten to twenty independent detections from one module at that cadence.
Raw transport bandwidth, settings and dropped frames must also be checked (\cref{app:hardware}).

Signal-to-noise ratio can inform a calibrated observation covariance, but a universal inverse-SNR weight ignores phase ambiguity, geometry and systematic alignment error.
A small mount-yaw error produces an approximately equal azimuth bias near the horizontal plane; a general three-dimensional pose error couples the angle channels.
Common calibration errors correlate observations across time and must not be averaged away as independent noise.

\subsection{Measurement Functions and Observability}

For a fixed sensor at world position $\vect s_i$, let $R_i$ map its calibrated local frame into the world frame.
Choose local axes consistently with the book's forward/up/right convention and explicitly convert any vendor-native axes.
For ball state $\vect x=(\vect p,\vect v)$, define
\begin{equation}
\vect q_i=R_i^{\mathsf T}(\vect p-\vect s_i),\qquad
\vect w_i=R_i^{\mathsf T}\vect v,\qquad \rho_i=\lVert\vect q_i\rVert.
\label{eq:implsensorframe}
\end{equation}
A candidate observation vector is
\begin{equation}
\label{eq:implmeasurements}
h_i(\vect x)=\begin{pmatrix}
\rho_i\\
\vect q_i^{\mathsf T}\vect w_i/\rho_i\\
\operatorname{atan2}(q_{iz},q_{ix})\\
\operatorname{atan2}\!\left(q_{iy},\sqrt{q_{ix}^2+q_{iz}^2}\right)
\end{pmatrix}.
\end{equation}
Use only the components actually supplied by the sensor and supported by its calibration; the OPS243-A speed channel does not supply all four.
Range rate requires $\rho_i>0$, and azimuth is undefined on the local vertical axis.
The elevation value approaches $\pm\pi/2$ there, but its Cartesian Jacobian is also singular; use a nonsingular local representation or exclude that geometry from an angle-based linearization.
Wrap angular innovations consistently and propagate pose and timestamp uncertainty.
Moving sensors require their own translational and rotational motion in the observation model.

Writing $\uvec u_i=(\vect p-\vect s_i)/\rho_i$ and $v_{ri}=\uvec u_i^{\mathsf T}\vect v$, the radial-speed Jacobian blocks are
\begin{equation}
\frac{\partial v_{ri}}{\partial\vect p}
=\frac{(\vect v-v_{ri}\uvec u_i)^{\mathsf T}}{\rho_i},
\qquad
\frac{\partial v_{ri}}{\partial\vect v}=\uvec u_i^{\mathsf T}.
\label{eq:implradialjacobian}
\end{equation}
Even with position known, a single instantaneous radial-speed observation leaves two velocity directions unconstrained.
Additional times, sensor viewpoints and dynamical assumptions can add information, but nearly repeated geometry can remain poorly conditioned.
Inspect local sensitivity over the actual observation window and compare posterior uncertainty with the prior; adding a state variable does not create a measurement of it.
In particular, including spin states does not establish that a short flight segment identifies spin magnitude, axis, wind and aerodynamic coefficients separately (\cref{sec:spinaxis}).

\subsection{Filtering, Smoothing, and the Impact Event}

An extended Kalman filter (EKF) and an extended Rauch--Tung--Striebel smoother are approximate nonlinear estimators, based on local linearization~\cite{sarkkasvensson2023}.
A proposed translational model is $\dot{\vect p}=\vect v$ and $\dot{\vect v}=\vect g+\vect a_{\rm aero}$, with a declared aerodynamic law and its parameter uncertainty.
If spin is omitted from the state, uncertain spin-dependent acceleration must still be represented through nuisance parameters or justified model-error covariance.
It cannot silently be treated as a known force.
Choose numerical time steps and interpolation by convergence checks at the required measurement precision.

Propagate between acquisition times, update only with associated valid observations, and retain the predicted and filtered means, covariances and transition Jacobians for backward smoothing.
For innovation $\vect\nu_k=\vect y_k-h_k(\widehat{\vect x}_k^-)$, a locally independent Gaussian observation model gives
\begin{equation}
S_k=H_kP_k^-H_k^{\mathsf T}+R_k,
\qquad d_k^2=\vect\nu_k^{\mathsf T}S_k^{-1}\vect\nu_k.
\label{eq:implinnovation}
\end{equation}
Here $H_k$ is the measurement Jacobian, $P_k^-$ the predicted state covariance and $R_k$ the observation-noise covariance; $R_k$ is distinct from the sensor rotation $R_i$ above.
A prespecified chi-square gate is only an approximate screening rule under the stated model.
Multipath or a misassociated club return can pass it, while a valid shot can fail a misspecified filter.
Inspect residual bias, temporal correlation, uncertainty coverage and rejection rates on held-out data.

Let a raw acoustic timestamp, expressed on the common clock, include sound travel over source-to-microphone distance $d$ at speed $c_s$ and electronic/detection delay $\tau_e$.
The estimated acoustic source-event time is
\begin{equation}
\widehat t_{\rm event}=t_{\rm stamp}-d/c_s-\tau_e.
\label{eq:impleventtime}
\end{equation}
The sound source event need not coincide exactly with face separation; establish that offset independently.
Estimate $d$ from the source location and microphone pose, propagate their uncertainty, and avoid subtracting delays already removed by acquisition software.
At $c_s=\SI{343}{\meter\per\second}$, one meter adds about \SI{2.92}{\milli\second}, much more than an assumed \SI{10}{\micro\second} electronic delay.
To first order, \SI{100}{\micro\second} of timing uncertainty at 150\,mph contributes about \SI{6.71}{\milli\meter} of along-track position uncertainty before other errors are included.
A post-impact flight smoother may extrapolate to the declared separation event, with increasing uncertainty outside its observed interval; it must not propagate a smooth free-flight model through the contact impulse into the pre-impact state.
Clock calibration and event definition are therefore part of the launch measurement, not merely data-logging details.

\subsection{Alignment Calibration Procedure}

Use surveyed targets or a calibration artifact to identify the sensor-to-world transforms, with several positions spanning the intended observation region.
Optical floor charts are useful precedents~\cite{uneekorcalib}, but radar needs a target and motion visible to its own measurement channels.
A ball rolled or swung along a surveyed line can test selected geometric relationships; rolling contact or a curved pendulum path is not automatically a known free-flight trajectory.
Validate range, direction, Doppler scale and clock alignment with separate checks where possible.
Record poses, calibration residuals, covariance assumptions, firmware, configuration and environmental conditions for each session.
Use independent poses or trajectories to test calibration transfer, rather than reporting the fitting residual as general accuracy.

\section{Derived Club Data, With Priors}
\label{app:priors}

The forward collision problem and inverse club-state problem have different information requirements (\cref{sec:dplane,sec:faceangle}).
For a forward calculation, declare contact-point incident velocity, local face normal, impact location, inertial properties and contact law.
The empirical relations in \cref{eq:tutelmanlaunch,eq:spinrate} are conditional approximations, not mutually independent physical identities~\cite{tutelman3d}.
Check that translation and spin remain compatible with the same impulse model and its domain of calibration.

For a local scalar launch approximation, write $b=wF+(1-w)P$, where $b$, $F$ and $P$ are launch, face-normal and path angles in one consistent projection.
For known $w\ne0$, its inverse and sensitivities are
\begin{align}
F&=P+\frac{b-P}{w},\label{eq:implfaceinverse}\\
\left(\frac{\partial F}{\partial b},\frac{\partial F}{\partial P},
\frac{\partial F}{\partial w}\right)
&=\left(\frac1w,-\frac{1-w}{w},-\frac{b-P}{w^2}\right).
\label{eq:implfacesensitivity}
\end{align}
Use these derivatives with the joint input covariance and a separate allowance for model discrepancy.
For illustration, $w=0.87$ gives a launch-angle sensitivity of about 1.15, but it is not a universal amplification factor for a three-dimensional mounting error.
Study-specific weights cannot be turned into a universal driver-to-wedge interpolation solely from dynamic loft (\cref{sec:facepathweighting}).
Separate horizontal and vertical inversions also require checking the validity of the small-angle projection approximation.

A club-type prior can regularize an underdetermined estimate, but the reported face angle, loft or spin then depends on that prior.
Record the prior distribution, its population and calibration range, and test sensitivity to plausible alternatives.
Do not use a spin fallback estimated from club priors as an independent confirmation of the same club estimate.
Unknown strike location, head inertia and local face geometry introduce gear-effect ambiguity; a fixed spin-axis allowance per millimeter of miss is not a universal uncertainty bound (\cref{sec:gear}).
Where those effects cannot be bounded from evidence, report the inverse quantity as unresolved or explicitly model-dependent.

\section{Optical Module (Phase 2) Design Parameters}
\label{app:optical}

The Wintriss disclosures and PiTrac provide optical concepts and implementation examples, respectively~\cite{us7292711,pitrac}.
They do not establish the accuracy or timing of a new camera, lens, illumination and trigger combination.
The following manufactured budget illustrates what must be checked.

\subsection{Shutter, Illumination, and Motion Blur}

The Raspberry Pi Global Shutter Camera uses an IMX296 sensor with $1456\times1088$ output, and its documented short-exposure capability is about \SI{30}{\micro\second} with adequate light~\cite{pigsproduct,pigscam}.
Global exposure avoids row-dependent capture times; it does not remove motion during integration.
The standard camera includes an infrared-cut filter, so an \SI{850}{\nano\meter} illuminator requires an explicitly verified optical configuration~\cite{pigscam}.
Do not assume unmodified spectral sensitivity or an arbitrary shorter exposure from the sensor name.

For locally transverse translation at speed $v_\perp$ and effective illumination duration $\tau$, object-plane blur is approximately $v_\perp\tau$.
At 150\,mph and \SI{30}{\micro\second}, this is \SI{2.01}{\milli\meter}, about 4.71\% of a nominal \SI{42.67}{\milli\meter} ball diameter.
Pixel blur is that fraction of the observed ball diameter in pixels, with additional effects from spin, defocus and lens distortion.
A shorter light pulse can determine the effective exposure only if it falls within the active integration interval and ambient light contributes negligibly.
Choose pulse width from the allowable image error and measured photon budget, not from a universal ``frozen motion'' threshold.

\subsection{Multiple Exposures and Trigger Timing}

For same-frame strobed images at approximately constant depth, pulse-center spacing $\Delta t$ produces center separation $v_\perp\Delta t$.
At 150\,mph, \SI{300}{\micro\second} and \SI{500}{\micro\second} give \SI{20.1}{\milli\meter} and \SI{33.5}{\milli\meter}, both less than the nominal ball diameter.
The silhouettes overlap; an algorithm must explicitly handle that overlap or use a different capture schedule.
An ideal no-overlap condition, before blur and margins, is
\begin{equation}
\Delta t>\frac{D}{v_\perp},\qquad
(N-1)v_\perp\Delta t+D\le L,
\label{eq:implstrobebudget}
\end{equation}
where $D$ is diameter, $N$ the number of pulses and $L$ the usable lateral field of view at the ball plane.
The first bound is about \SI{636}{\micro\second} for this example.
For a rectangular light pulse of duration $\tau$ with pure transverse translation, replace $D$ by $D+v_\perp\tau$ in both bounds to include the illuminated motion extent; additional errors still require margins.
Five pulses spaced by \SI{0.8}{\milli\second} occupy about \SI{257}{\milli\meter}, including one ball diameter, within an ideal \SI{300}{\milli\meter} field before allowance for blur, pose, depth change and timing error.
These numbers specify a test case, not a commercial capture geometry or a validated design.

Budget acoustic propagation, threshold-crossing jitter, transport delay and camera/strobe synchronization together with the event offset in \cref{eq:impleventtime}.
Use pretrigger buffering where appropriate; an asserted microphone-module delay alone does not solve impact timing.
Verify actual pulse times and integration windows electrically and with a moving target.

\subsection{Position, Spin, and Club Pose}

Circle detection is a candidate image-processing step, not proof that partially overlapping, blurred or occluded balls are correctly localized.
A known-diameter depth estimate requires calibrated projection, actual ball size and a suitable silhouette model; stereo additionally requires calibrated extrinsics, correspondence and synchronization.
For close views, account for the perspective silhouette of a sphere rather than assuming the far-field diameter formula is exact.

Texture registration across ball images can estimate rotation on $SO(3)$ when visible features constrain it~\cite{us7292711,foresightspherical}.
At 3,000\,rpm, \SI{400}{\micro\second} corresponds to $7.2\degs$, but a small angle alone proves neither uniqueness nor convergence of a local search.
Dimple repetition, partial visibility, lighting and image noise can leave multiple compatible rotations.
Use multiple hypotheses or intervals, inspect residuals and determine which rotation components are constrained before converting an increment to a spin vector.
A D-plane prior may guide a search; it must not force the answer and then be counted as independent optical evidence.
Calibrated face fiducials or additional viewpoints can constrain club pose and strike location, with their own geometry and uncertainty requirements~\cite{us5501463,us6758759}.

\section{Flight Model and Validation Protocol}

Treat aerodynamic coefficients, spin decay, wind, ball properties and numerical integration as parts of a declared model, not universal constants inherited from a named publication.
A literature model can supply a starting hypothesis~\cite{smitssmith}; trajectory fitting can estimate parameters~\cite{us6186002}, but neither guarantees identification over a short observation arc.
Carry alone is one scalar outcome and generally cannot determine a complete lift, drag and spin-decay law.
Version the model and coefficient domain, verify signs and units, and test integration convergence and event detection before interpreting a physical discrepancy.
Do not infer golf accuracy from solver convergence alone.

Separate calibration trajectories from held-out validation by the relevant independent units, such as session, ball batch or participant.
Prespecify reference definitions, practical error limits, exclusions and missing-shot reporting (\cref{ch:accuracy,ch:validation-program}).
Paired commercial measurements test agreement, including the reference instrument's uncertainty; they are not automatically ground truth.
Use repeated-measures analysis when shots share a participant or session, and examine condition-dependent bias and spread rather than relying on correlation or one pooled mean~\cite{blandaltman2007}.

For MLM2PRO spin comparison, Rapsodo specifies RPT printed balls; an RCT radar ball is not an interchangeable optical reference~\cite{rapsodomlm2prorpt}.
Record ball identity, device mode, firmware, failed reads and whether every reported channel was measured, derived or estimated.
Log raw I/Q, available frames, timestamps and calibration/configuration identifiers so an algorithm revision can be replayed, while retaining an untouched test set for final evaluation.
Replayability makes errors inspectable; independent observations and appropriate uncertainty analysis establish what the system actually measures.

% Keep the existing background catalog entries; they do not substantiate the
% removed universal launch-weight, smash-ceiling or deployed-algorithm claims.
\nocite{trackmanballflightlaws,trackmanspinpatentblog,tutelmansmash}
